\section{Multiplying by a character before turning}\label{cs:characters}

For an input class \(E\), define
\begin{equation}\label{cs:decorated}
 D_j^\epsilon[E]=\Turn(R^\epsilon L_jE).
\end{equation}
The character stays at its specified position before the turn. The
input may be a product of several subtree polynomials; it need not
have a positive coefficient expansion.

\begin{lemma}[Coefficient duality]\label{cs:duality}
For \(0\le r\le d\),
\begin{equation}\label{cs:dual}
 \coef{r}{D_j^\epsilon[E]}=\phi(e_rR^\epsilon L_jE),\qquad
 |\coef{r}{D_j^\epsilon[E]}|\le CN^A\normJ{E}w_{N,r}.
\end{equation}
\end{lemma}
\begin{proof}
The pairing in \eqref{cs:turn-trace} extracts the coefficient
\(d-r\) before the turn, hence coefficient \(r\) after the turn.
Use \eqref{cs:finite-cost}, submultiplicativity and
\eqref{cs:basis-norm}. There is no sum over the coefficients of
the input representative.
\end{proof}

Set \(G_0=g(1/4)\). Differentiation gives
\begin{equation}\label{cs:g-derivatives}
 g'(u)=2\log\cot(\pi u),\qquad
 g''(u)=-\frac{4\pi}{\sin(2\pi u)}<0,\qquad 0<u<1/2.
\end{equation}
Thus \(G_0>0\) is the unique maximum of \(g\), while
\(g(0)=g(1/2)=0\).

\begin{lemma}[Angular bound for a low Newton polynomial]\label{cs:low-strip}
Let \(I\) be a compact subinterval of \((0,1/2)\). There are
\(c_I,C_I,y_I>0\) and \(e_N\to0\) such that, whenever \(r/N\in I\),
\begin{equation}\label{cs:low-gaussian}
 |e_r(2\cos(\theta+\ii y))|
 \le w_{N,r}\exp N\{-c_I(\theta-\pi)^2+C_Iy^2+e_N\}
\end{equation}
for \(0\le\theta\le2\pi\), \(|y|\le y_I\).
On the full low range, the upper logarithmic rate is
\begin{equation}\label{cs:Newton-action}
 b(u,\theta,y)=\int_0^u
      \log|2\cos(2\pi v)-2\cos(\theta+\ii y)|\dd v-2J(u).
\end{equation}
At \(y=0\), \(g(u)+b(u,\theta,0)\) has the unique maximum
\(2G_0\) at \((u,\theta)=(1/4,\pi)\).
\end{lemma}
\begin{proof}
The numerator product in \eqref{cs:basis} is a midpoint Riemann
sum for the integral in \eqref{cs:Newton-action}. Factor each cosine
difference into its two sine factors. For an upper bound,
replace each of their logarithms separately by its maximum with \(-K\). The resulting
continuous functions have uniform Riemann sums on the compact parameter
set. The negative tails of the two circular logarithms have integrals
bounded uniformly by \(C e^{-K}(1+K)\): near a zero this follows by
integrating \(|\log|t||\) on a short interval, and moving the zero off
the real circle only decreases its negative tail. This also covers
coalescing zeros. Let \(K\to\infty\) after taking the upper limit in
\(N\). The factorial estimate in equation~\eqref{cs:factorial} handles
the denominator, including both endpoints. This proves a uniform
upper error \(o(1)\) in the rate.

Here is the strict real comparison. Write
\(\ell(t)=\log|2\sin(t/2)|\), a periodic function of period \(2\pi\).
For \(0<u<1/2\), the numerator integral is
\[
 \frac1{2\pi}\int_{\theta-2\pi u}^{\theta+2\pi u}\ell(t)\dd t.
\]
As the center \(\theta\) moves from \(0\) to \(\pi\), this integral
strictly increases. Its derivative is the difference of its endpoint
values divided by \(2\pi\); that difference is positive because
\[
 \sin^2\frac{\theta+\alpha}{2}
       -\sin^2\frac{\theta-\alpha}{2}
       =\sin\theta\sin\alpha>0,
 \qquad \alpha=2\pi u\in(0,\pi).
\]
The isolated logarithmic singularities are integrable, so the argument
extends across them by continuity. Reflection treats \(\pi<\theta<2\pi\).
It follows that
\(b(u,\theta,0)\le b(u,\pi,0)=g(u)\), strictly away from \(\pi\)
when \(0<u<1/2\). At \(u=0,1/2\) the rate is zero. Equation
\eqref{cs:g-derivatives} proves the stated unique joint maximum.

Near \(\theta=\pi\), put \(w=\theta-\pi+\ii y\). The numerator
has no zeros for \(u\in I\) in a small fixed neighborhood, and its
holomorphic action there is
\begin{equation}\label{cs:B-action}
 B(u,w)=-J(1/2-u+w/(2\pi))-J(1/2-u-w/(2\pi))-2J(u).
\end{equation}
The analytic branch of \(J\) is the real one on this compact interior
interval. It satisfies
\[
 B(u,0)=g(u),\quad B_w(u,0)=0,\quad
 B_{ww}(u,0)=-\frac{\tan\pi u}{2\pi}<0.
\]
Uniform Taylor bounds give a negative angular quadratic and a positive
quadratic height cost there. The remaining real compact set has a
strict gap by the preceding comparison. Continuity of the logarithmic
integrals preserves half that gap at sufficiently small fixed heights.
Reducing the angular constant gives \eqref{cs:low-gaussian} on the
whole strip. Finally use \eqref{cs:weight-rate} to replace
\(e^{Ng(r/N)}\) by \(w_{N,r}\).
\end{proof}

\begin{proposition}[Norm and angular localization after a character]\label{cs:decorated-turn}
Suppose \(\normJ{E_N}\le e^{N(c+o(1))}\), uniformly in the specified
input family. Then
\begin{equation}\label{cs:D-norm}
 \normJ{D_j^\epsilon[E_N]}\le e^{N(c+2G_0+o(1))}.
\end{equation}
Its canonical representative has constants \(a,C,y_0>0\) such that
\begin{equation}\label{cs:D-strip}
 |D_j^\epsilon[E_N](2\cos(\theta+\ii y))|
 \le e^{N(c+2G_0-a(\theta-\pi)^2+Cy^2+o(1))}
\end{equation}
on \(0\le\theta\le2\pi\), \(|y|\le y_0\). The constants are
independent of the physical color \(j\) and retained sign.
A fixed loss in the norm exponent of \(E_N\) is preserved in the
coefficient estimate \eqref{cs:dual} at every output index.
\end{proposition}
\begin{proof}
For low indices, reconstruct the polynomial using
\eqref{cs:dual} and \eqref{cs:basis-norm}. The norm is bounded by
a polynomial times
\(\normJ{E_N}\sum_{r\le M}w_{N,r}^2\), whose rate is at most
\(c+2G_0\). The high part has rate at most \(c\), proving
\eqref{cs:D-norm}.

For the low angular estimate, the rate of the \(r\)-th term is
\(c+g(r/N)+b(r/N,\theta,y)\). Lemma~\ref{cs:low-strip} gives
a strict index and angular quadratic at its unique maximum and a
strict gap on the compact complement. The same argument with a small
fixed height gives the required bound for this entire low sum. The
high part is itself canonical, has norm at most \(e^{N(c+o(1))}\),
and by \eqref{cs:ellipse} has strip rate at most \(c+|y|\).
Choose \(y_0\) and then the angular constant small enough that
\(y_0+a\pi^2<2G_0\). This part then lies below
\eqref{cs:D-strip} with a fixed loss. All sums have at most \(N\)
terms. The final assertion follows directly from \eqref{cs:dual}.
\end{proof}

\begin{proposition}[A deficit in the new character]\label{cs:regular-deficit}
Suppose an actual representative of \(E_N\) satisfies
\begin{equation}\label{cs:E-strip}
 |E_N(2\cos(\theta+\ii y))|
 \le e^{N(c-c_E(\theta-\pi)^2+C_Ey^2+o(1))},
 \quad |y|\le y_E,
\end{equation}
with \(c_E\ge0\). The empty input \(E_N=1,c=0\) is allowed.
For every compact \(I\subset(0,1/2)\) and every \(\delta>0\),
there is \(\eta>0\) such that
\begin{equation}\label{cs:regular-color-loss}
 N-j\ge\delta N,\quad r/N\in I
 \quad\Longrightarrow\quad
 |\coef{r}{D_j^\epsilon[E_N]}|
          \le e^{N(c+g(r/N)-\eta)}
\end{equation}
for all sufficiently large \(N\). If also
\(\normJ{E_N}\le e^{N(c+o(1))}\), then this subfamily has
\(\normJ{D_j^\epsilon[E_N]}\le e^{N(c+2G_0-\eta')}\)
for some \(\eta'>0\).
\end{proposition}
\begin{proof}
By Lemma~\ref{cs:low-strip}, the actual product \(e_rE_N\)
satisfies the positive-angular hypothesis of Theorem~\ref{cs:root-bound}
uniformly for \(r/N\in I\), at rate \(c+g(r/N)\).
The first identity in \eqref{cs:dual} and that theorem prove
\eqref{cs:regular-color-loss}. No continuation of the discrete color
\(j\) is used.

For the norm assertion take \(I\) to be a small interval about
\(1/4\). Its contribution loses \(\eta\). The other low indices
have a fixed loss from \(2g(u)<2G_0\) outside \(I\), and the high
part has the fixed gap \(2G_0\). Apply \eqref{cs:dual} to those
remaining terms and sum. Reducing the minimum loss proves the claim.
\end{proof}

For a later component, the relevant input index can be strictly above
\(1/4\). This is why \eqref{cs:regular-color-loss} is stated on every
compact regular interval, rather than only near the maximum of \(g\).
